\chapter{Conclusions}\label{ch:conclusions}

\section{What the evaluations establish}

\begin{itemize}[leftmargin=*]
  \item \textbf{The learned intent model is the copilot's main gain.} It beats keyword routing on held-out phrasings
  ($p = \est{int.mcnemar}{p}$) and on whole conversations ($p = \est{chl.mcnemar}{p}$), and its confidence is
  calibrated well enough for a threshold to be trusted: above it, \est{int.acc-confident}{pct}\,\% of held-out
  phrasings are labelled correctly.
  \item \textbf{Safety comes from the architecture, not the model.} Alone, models of every size and provider broke
  card-service rules; inside BETA AID none did, on any attack, and the refusals did not depend on the model. The model
  never holds what an attacker wants, never acts, and never speaks unchecked.
  \item \textbf{The language models do not separate on accuracy.} With a model inside, the copilot ends within a few
  cases of the copilot without one, and no difference is detectable at this sample size. Every model rescues the same
  informal phrasings the intent model is unsure of, and loses a few multi-turn or out-of-scope cases.
  \item \textbf{They do separate on speed and reliability.} Calls range from a steady sub-second median to timeouts
  and malformed replies, and the same model differs between providers.
  \item \textbf{Governed retrieval keeps unapproved documents out, but its threshold does not reject unanswerable
  questions.} The steps after retrieval keep that from reaching a customer as an unverified answer.
\end{itemize}

\section{Recommendation}

Run the copilot with \model{Qwen3-30B-A3B-Instruct-2507} on Nebius Token Factory. Inside BETA AID it was as safe as
every other model (\est{atk.qwen30-tf-wrapped}{num} violations in \est{atk.qwen30-tf-wrapped}{den} attacks) and as
accurate as the evaluations can tell (\est{mod.qwen30-tf-correct-compared}{num} of
\est{mod.qwen30-tf-correct-compared}{den} compared cases). The choice rests on the screening: the lowest median call
of the hosted options (\est{scr.qwen30-tf-p50}{estimate}\,ms), a slowest call well inside the timeout
(\est{scr.qwen30-tf-max}{value}\,ms), and a valid label every time (\est{scr.qwen30-tf-valid}{num} of
\est{scr.qwen30-tf-valid}{den}). The copilot remains fully usable without any model, with the same safety; the model
adds the informal phrasings it would otherwise ask the customer to rephrase.

\section{Limits and next steps}

\begin{itemize}[leftmargin=*]
  \item \textbf{Sample sizes.} The frozen set cannot rank models within a few cases of each other. A larger set,
  written the same way and frozen before scoring, is the way to separate them; repeated runs would also measure how
  much a model's answers vary from run to run.
  \item \textbf{Retrieval rejection.} The relevance threshold rests on a handful of unanswerable development
  questions; more of them, and a threshold that uses the gap between the top scores, would make the knowledge base say
  \enquote{no document} when it should.
  \item \textbf{The losses with a model.} The multi-card follow-ups and the out-of-scope request that models fail should
  be replayed with full traces, as the knowledge cases were, to locate the stage that produced them.
  \item \textbf{The defect in the stored runs.} The model runs of the local and gateway models should be repeated with
  retrieval working; only the hosted models' knowledge cases were replayed here.
  \item \textbf{Cost.} The runs recorded no prices. The tokens per call in \cref{ch:screening} price the job at any
  provider's rates, and recording prices with each run would put cost per resolved case beside accuracy.
\end{itemize}
